\chapter{Implementation / Construction}\label{ch:impl}

This chapter describes how the design of \Cref{ch:arch,ch:design} was built: the
environment the team worked in, the main modules with short excerpts of the real code,
the way the AI models are wired in, and the structure and standards of the code base.
Every listing is copied from commit \texttt{60d40f9}, with its file path and first line
number; long lines were kept as they are, and non-ASCII characters were replaced.

\section{Development Environment and Tools}\label{sec:devenv}

\Cref{tab:devenv} lists the environment used to build RecaLL.

{\small
\begin{longtable}{@{}p{0.22\textwidth}p{0.72\textwidth}@{}}
\caption{Development environment and tools}\label{tab:devenv}\\
\toprule
\textbf{Item} & \textbf{Details}\\
\midrule
\endfirsthead
\toprule
\textbf{Item} & \textbf{Details}\\
\midrule
\endhead
\bottomrule
\endfoot
Operating system & Windows on the team's four laptops (\Cref{tab:hw}); Ubuntu (GitHub-hosted) for CI; \texttt{windows-latest} for release builds.\\
Editors and IDEs & Antigravity and VS Code. The Antigravity agent generated most of the code from the requirements and designs the team had prepared, and GitHub Copilot opened two of the pull requests.\\
Backend runtime & Python 3.13.5 in a local virtual environment (\texttt{.venv}); Python 3.12 in CI and in the Docker image. Dependencies in \texttt{requirements.txt} (unpinned except \texttt{json-repair==0.63.3}), installed with \texttt{pip}.\\
Frontend runtime & Node.js 22 locally, Node 20 in CI; pnpm (lockfile \texttt{desktop/pnpm-lock.yaml}); Vite dev server on port 5173.\\
Desktop build & Rust stable toolchain and the Tauri CLI 2.11; PyInstaller for the backend sidecar (\texttt{scripts/build\_sidecar.py}); \texttt{build.bat} for local builds.\\
Local start-up & \texttt{run.bat} opens two windows: \texttt{uvicorn app.main:app --host 127.0.0.1 --port 8000 --reload} for the backend, and \texttt{pnpm dev --open} for the frontend. \texttt{stop.bat} stops both.\\
Configuration & \texttt{.env} (ignored by Git) holds the environment, secret keys, the default provider and optional server keys; \texttt{.env.example} documents the variables. Without values, the backend creates local secret files in the user data folder.\\
AI services for development & The Gemini API (configured model \texttt{gemini-3.6-flash}) and local models through Ollama; the other providers are options in the application. Langfuse optional.\\
Containers & \texttt{Dockerfile} (Python 3.12-slim, Uvicorn with 2 workers, health check on \texttt{/health}), \texttt{desktop/Dockerfile} (Node 20 build, nginx 1.27 runtime) and \texttt{docker-compose.yml} (backend, frontend, PostgreSQL and Ollama). These support a hosted demo; the product itself is the desktop installer.\\
Version control & Git and GitHub (\Cref{sec:git}); GitHub Actions for CI and releases (\Cref{sec:cicd}).\\
Database tools & Schema documentation generated with tbls into \texttt{docs/database/}.\\
Testing tools & pytest 9.2, FastAPI \texttt{TestClient}, Playwright 1.62 (Chromium), oxlint 1.74, \texttt{tsc}.\\
\end{longtable}
}

\section{Key Implementation Details / Modules Developed}\label{sec:modules}

\subsection*{Document processing}
A PDF imported from the interface goes through the study path:
\texttt{POST /books/upload} stores the file, reads the PDF's table of contents (falling
back to heading detection in \texttt{toc\_parser.py}), filters front matter and builds
the topic tree. Topic text is then extracted on demand through the extractor registry
(\texttt{pdf\_extract.py}: PyMuPDF4LLM by default, Marker for scanned pages when it is
installed), cleaned and segmented into atomic concepts by the LLM. The same books are
mirrored into \texttt{recall\_saas.db} by \texttt{BookSyncService}, and documents
uploaded through \texttt{/api/v1/documents/upload} go through
\texttt{DocumentProcessingPipeline}. \Cref{lst:pipeline} shows the chunking and
embedding stages of that pipeline: each stage first records its name and progress in the
job row, and an empty document still produces one fallback chunk so that later stages
never see an empty list.

\begin{lstlisting}[language=Python,caption={Chunking and embedding stages (\texttt{app/services/pipeline.py}, lines 137--171; lines 156--166 elided)},label={lst:pipeline}]
            ProcessingJobRepository.update_stage(
                job_id=job_id,
                current_stage="chunking",
                stage_progress=60,
                status="running"
            )

            chunker = SemanticChunker(target_chunk_size=1500, overlap_size=150)
            chunks = chunker.chunk_document(extracted.pages)

            if not chunks:
                # Handle empty or zero-text documents with a fallback single chunk
                chunks = [
                    SemanticChunker.chunk_text(
                        text=extracted.full_text or "No extractable text content found.",
                        page_number=1,
                        start_index=0
                    )[0]
                ]
            # ... (stage marker for "embedding", progress 80)
            chunk_texts = [c.content for c in chunks]
            embeddings = EmbeddingService.generate_embeddings(
                texts=chunk_texts,
                workspace_id=workspace_id
            )
\end{lstlisting}

\subsection*{Cornell notes}
\texttt{generate\_single\_cornell\_note()} in \texttt{app/main.py} waits for any running
extraction of the topic, makes sure the topic has Markdown, lists up to 12 atomic
concepts and every figure found in the text, and fills the
\texttt{cornell\_scaffold\_prompt} (\Cref{lst:cornell}). The prompt asks for a fixed
structure: a core intuition, invariants and definitions, an analogy, a first-principles
breakdown, equation anatomy in LaTeX, worked examples, exam pitfalls and self-testing cue
questions~\cite{pauk2011}. For a chapter, the notes of the sub-topics are
assembled into one master guide, and the cue questions are collected at the end. The
\texttt{scaffold-stream} route sends progress events while this runs.

\begin{lstlisting}[language=Python,caption={Filling the Cornell prompt (\texttt{app/main.py}, from line 2347)},label={lst:cornell}]
    from app.prompt_manager import get_prompt_template
    scaffold_prompt = get_prompt_template("cornell_scaffold_prompt").format(
        topic_title=topic.get('title', 'Study Topic'),
        breadcrumb=topic.get('breadcrumb', ''),
        concepts_summary=concepts_summary,
        content_reference=content[:8000],
        diagrams_reference=diagrams_reference,
    )

    local_settings = copy(settings)
    provider = get_llm_provider(local_settings, provider_override=provider_override)
    scaffold_markdown = await provider.generate(scaffold_prompt, json_schema=None, temperature=0.2, max_tokens=4096, feature="cornell_notes")
    if not scaffold_markdown or not scaffold_markdown.strip():
        raise RuntimeError(f"LLM provider returned empty response for study notes of '{topic.get('title')}'.")
\end{lstlisting}

\subsection*{Socratic drill}
The drill has two calls. \texttt{generate\_diagnostic\_questions()} asks for exactly two
questions: the first must test a causal mechanism, the second a counterfactual or an
applied scenario, each with 2--4 key invariants and a hint (temperature 0.3).
\texttt{evaluate\_student\_answer()} grades the answer with the
\texttt{socratic\_evaluation\_prompt} at temperature 0.1 (\Cref{alg:socratic}). Both use
\texttt{\_generate\_and\_parse()} (\Cref{lst:drill}): a provider error is logged and the
call is tried again, and an unparseable answer is retried once with a stricter
instruction appended to the prompt. The parsed object is then validated by the
\texttt{DiagnosticQuestionSet} or \texttt{DiagnosticEvaluation} Pydantic model.

\begin{lstlisting}[language=Python,caption={Generate, repair and retry (\texttt{app/socratic\_drill.py}, from line 389; an em dash in line 413 replaced by \texttt{--})},label={lst:drill}]
    for attempt in range(1, max_attempts + 1):
        try:
            raw = await provider.generate(
                prompt=current_prompt,
                json_schema=schema,
                temperature=temperature,
                max_tokens=max_tokens,
                feature="socratic_drill",
            )
        except RecallError:
            raise
        except Exception as e:
            logger.error(f"{context}: LLM provider call failed on attempt {attempt}/{max_attempts}: {e}")
            last_error = e
            continue

        try:
            return _parse_llm_json(raw, context=context)
        except ValueError as e:
            logger.warning(f"{context}: attempt {attempt}/{max_attempts} produced unparseable JSON: {e}")
            last_error = e
            current_prompt = (
                prompt
                + "\n\nIMPORTANT: Your previous response was not valid JSON "
                  f"({e}). Respond ONLY with a single valid JSON object -- "
\end{lstlisting}

\subsection*{FSRS review}
The review screens call the \texttt{/api/v1/reviews} routes. A session takes up to 20
due items ordered overdue $\rightarrow$ due now $\rightarrow$ new; each item is shown
front first, revealed, and rated. \Cref{lst:fsrs} is the core of
\texttt{FSRSScheduler.schedule()}: the card stored as JSON is restored (or created), the
\texttt{fsrs} library computes the new card, and the due time and interval are derived
from it~\cite{ye2022fsrs}. \texttt{ReviewService.rate\_item()} then writes the learning
item, a review event, the session item and the session counters in one transaction.

\begin{lstlisting}[language=Python,caption={Applying an FSRS rating (\texttt{app/services/scheduler.py}, from line 159)},label={lst:fsrs}]
        # 3. Restore or initialize FSRS Card
        meta = dict(current_metadata or {})
        card_data = meta.get("fsrs_card")
        if card_data and isinstance(card_data, dict):
            try:
                card = FSRSCard.from_dict(card_data)
            except Exception as e:
                logger.warning("Failed to restore FSRS card from dict: %s. Reinitializing.", e)
                card = FSRSCard()
        else:
            card = FSRSCard()

        # 4. Apply FSRS review calculation
        updated_card, review_log = self._fsrs.review_card(card, fsrs_rating, eval_time)

        # 5. Extract next review timestamp and interval
        next_review_at = updated_card.due
        if next_review_at.tzinfo is None:
            next_review_at = next_review_at.replace(tzinfo=timezone.utc)

        interval_seconds = max(0.0, (next_review_at - eval_time).total_seconds())
        scheduled_days = round(interval_seconds / 86400.0, 2)
\end{lstlisting}

\subsection*{Knowledge Hub (RAG)}
The Knowledge Hub view lists conversations and lets the student limit a conversation
to one document. \texttt{RAGService.execute\_rag()} checks that the conversation belongs
to the workspace, reserves one credit unless a BYOK key is used, saves the question,
retrieves the five best chunks (\Cref{alg:retrieval}), builds the prompt with the
untrusted text inside \texttt{<reference\_data>} tags and the last six messages, calls
the provider, validates citations and saves the answer (\Cref{lst:rag})~\cite{lewis2020rag}.
A streaming variant (\texttt{execute\_rag\_stream}) follows the same steps.

\begin{lstlisting}[language=Python,caption={Retrieval, generation and citation check (\texttt{app/services/rag.py}, from line 259)},label={lst:rag}]
            # 4. Hybrid Knowledge Retrieval
            search_res = RetrievalService.search(
                query=query,
                workspace_id=workspace_id,
                document_id=document_id,
                mode="hybrid",
                limit=5
            )
            retrieved_items = search_res.results

            # 5. Fetch history & build isolated prompt
            history = MessageRepository.list_by_conversation(conversation_id, limit=cls.MAX_HISTORY_MESSAGES + 1)
            # Exclude the user message just inserted
            prior_history = [m for m in history if m["role"] == "assistant" or (m["role"] == "user" and m["content"] != query)]
            ai_messages = cls._build_context_prompt(query, retrieved_items, prior_history)

            # 6. Invoke AI Provider
            try:
                ai_resp, provider_name, model_name, is_byok = await AIService.generate(
                    workspace_id=workspace_id,
                    messages=ai_messages,
                    provider_preference=provider_preference,
                    temperature=temperature,
                    max_tokens=max_tokens
                )
\end{lstlisting}

\subsection*{Quizzes}
\texttt{QuizGenerationService.generate\_quiz()} retrieves chunks from the selected documents
(spread evenly across them), asks the model for multiple-choice and true/false
questions with explanations and source references, and keeps only questions that pass
the checks of \Cref{sec:ai_design}. Grading is deterministic and never uses the model
(\Cref{lst:quiz}): an unanswered question counts as wrong, multiple-choice answers are
compared by index or by text, and true/false answers accept several spellings. The
result updates the learning items, so missed questions enter the review queue.

\begin{lstlisting}[language=Python,caption={Deterministic grading (\texttt{app/services/attempts.py}, from line 242)},label={lst:quiz}]
        for q in questions:
            q_id = q["id"]
            q_type = q["type"]
            stored_correct = str(q["correct_answer"]).strip()
            selected = answers_map.get(q_id)

            if selected is None or selected == "":
                # Unanswered: counted as incorrect
                is_correct = False
                selected_str = ""
            else:
                selected_str = str(selected).strip()
                if q_type == "multiple_choice":
                    # Evaluate index equality or text equality
                    if selected_str.isdigit() and stored_correct.isdigit():
                        is_correct = (int(selected_str) == int(stored_correct))
                    else:
                        is_correct = (selected_str.lower() == stored_correct.lower())
                elif q_type == "true_false":
                    sel_bool = selected_str.lower() in ("true", "t", "1", "yes")
                    correct_bool = stored_correct.lower() in ("true", "t", "1", "yes")
\end{lstlisting}

\subsection*{Analytics}
The dashboard and Progress views read \texttt{/api/v1/learning/*}. The SQL aggregation
is in \texttt{AnalyticsRepository}; \texttt{LearningAnalyticsService} turns it into
response models. \Cref{lst:analytics} shows how the daily activity series is zero-filled,
so that the Recharts charts have one point per day even on days without study.

\begin{lstlisting}[language=Python,caption={Zero-filled daily activity (\texttt{app/services/learning\_analytics.py}, from line 103)},label={lst:analytics}]
        for d_str in date_keys:
            act = raw_activity.get(d_str, {
                "reviews_count": 0,
                "quiz_attempts_count": 0,
                "correct_answers": 0,
                "incorrect_answers": 0
            })
            r_cnt = act.get("reviews_count", 0)
            q_cnt = act.get("quiz_attempts_count", 0)
            c_cnt = act.get("correct_answers", 0)
            i_cnt = act.get("incorrect_answers", 0)

            if r_cnt > 0 or q_cnt > 0:
                total_active_days += 1
            total_reviews += r_cnt
            total_quizzes += q_cnt
\end{lstlisting}

\subsection*{AI providers and BYOK}
In Settings, each provider card tests the key against the provider and saves it through
\texttt{POST /api/v1/account/byok/credentials}. The key is encrypted with AES-256-GCM
under a 32-byte master key (\texttt{BYOK\_ENCRYPTION\_KEY}, or a generated local key
file), and only a masked hint is returned (\Cref{lst:byok}). When a feature runs,
\texttt{AIService.resolve\_adapter()} decrypts the key and builds the adapter.

\begin{lstlisting}[language=Python,caption={Encrypting a BYOK key (\texttt{app/core/security.py}, from line 135)},label={lst:byok}]
def encrypt_secret(plaintext: str, key_override: Optional[bytes] = None) -> Tuple[str, str, str]:
    """
    Encrypts a secret (e.g., BYOK API key) using AES-256-GCM.
    Returns:
        (ciphertext_hex, nonce_hex, tag_hex)
    """
    key = key_override or get_encryption_master_key()
    aesgcm = AESGCM(key)
    nonce = secrets.token_bytes(12)  # 12-byte GCM standard nonce

    # AESGCM.encrypt returns ciphertext + 16-byte tag appended
    encrypted_data = aesgcm.encrypt(nonce, plaintext.encode('utf-8'), None)
    ciphertext = encrypted_data[:-16]
    tag = encrypted_data[-16:]

    return ciphertext.hex(), nonce.hex(), tag.hex()
\end{lstlisting}

\section{AI-integration Implementation}\label{sec:ai_impl}

No model is trained, so there are no training scripts, datasets or model files; the
integration code is the HTTP clients, the prompt registry and the output handling.

\paragraph{Pre-processing.} Extracted Markdown is cleaned before it reaches a prompt:
\texttt{markdown\_cleanup.py} can repair extracted Markdown with the
\texttt{markdown\_cleanup\_prompt}, \texttt{latex\_fixer.py} wraps bare LaTeX left by Marker in math
delimiters, image tags are removed for the drill, and every input is cut to its character
or token budget (\Cref{sec:ai_design}).

\paragraph{Model loading and inference.} Hosted models need no loading; the providers
build one \texttt{httpx.AsyncClient} per request with a 120-second time-out (study stack)
and send the prompt, the temperature, the token limit and, for structured tasks, the JSON
schema. For Ollama, the user runs the model server; RecaLL only checks that it is
reachable (\texttt{/settings/verify-ollama}) and calls its HTTP API. Streaming answers
(tutor chat, Knowledge Hub) are forwarded to the interface as SSE.

\paragraph{Retries and rate limiting.} \Cref{lst:retry} shows the core of
\texttt{retry\_with\_backoff()}: non-retriable status codes are raised at once, the
server's \texttt{Retry-After} header is honoured, and otherwise the delay grows by a
factor of two with random jitter. The Gemini provider calls
\texttt{rate\_limiter.acquire("gemini", priority)} first; the limiter keeps a sorted
queue of waiters per provider, so an interactive request (priority 0) is served before
queued batch requests (priority 1).

\begin{lstlisting}[language=Python,caption={Exponential back-off with jitter (\texttt{app/llm\_providers/retry.py}, from line 50)},label={lst:retry}]
    for attempt in range(max_retries + 1):
        try:
            return await coro_factory()
        except httpx.HTTPStatusError as e:
            status = e.response.status_code

            # Non-retriable: fail fast on client errors (Auth, Bad Schema, Forbidden, etc.)
            if status not in RETRIABLE_STATUS_CODES or attempt == max_retries:
                logger.warning(
                    f"[{provider_name}] Non-retriable or exhausted HTTP {status} on attempt {attempt + 1}/{max_retries + 1}: {e}"
                )
                raise

            # Check if provider returned a Retry-After header
            retry_after_hdr = e.response.headers.get("retry-after") or e.response.headers.get("Retry-After")
            if retry_after_hdr:
                try:
                    sleep_time = min(float(retry_after_hdr), max_delay)
                except (ValueError, TypeError):
                    sleep_time = min(delay + random.uniform(0, jitter), max_delay)
            else:
                sleep_time = min(delay + random.uniform(0, jitter), max_delay)
                delay *= backoff_factor
\end{lstlisting}

\paragraph{Structured output.} The v1 services parse model output with
\texttt{parse\_structured\_json()} (\Cref{lst:json}), which strips code fences and tries
plain JSON first, then \texttt{json\_repair} on the outermost object or array, then on
the whole text~\cite{json_repair}. The study stack uses the stricter
\texttt{\_parse\_llm\_json()} pipeline described in \Cref{sec:ai_design}.

\begin{lstlisting}[language=Python,caption={Tolerant JSON parsing (\texttt{app/services/ai/base.py}, from line 117)},label={lst:json}]
    if not text or not text.strip():
        return None

    cleaned = text.strip()

    # 1. Strip Markdown code fences if wrapped
    if cleaned.startswith("```"):
        cleaned = re.sub(r"^```(?:json)?\s*", "", cleaned, flags=re.IGNORECASE)
        cleaned = re.sub(r"\s*```\s*$", "", cleaned)
        cleaned = cleaned.strip()

    # 2. Try standard json.loads if it yields a container (dict or list)
    try:
        val = json.loads(cleaned)
        if isinstance(val, (dict, list)):
            return val
    except Exception:
        pass

    # 3. Extract outer JSON object or array bounds
    first_brace = cleaned.find("{")
    first_bracket = cleaned.find("[")
\end{lstlisting}

\paragraph{Prompt templates and overrides.} \texttt{get\_prompt\_template(key)} returns
the user's saved override from the \texttt{system\_prompts} table if there is one, and
the registry default otherwise, wrapped in \texttt{SafePromptTemplate} so that a missing
optional variable becomes an empty string instead of an exception.

\paragraph{Caching, cost and latency.} Generated artefacts are stored and reused rather
than regenerated: topic Markdown, concepts, Cornell notes, flashcards, summaries
(\texttt{GET /documents/\{id\}/summary} makes no AI call) and chat history. Duplicate
flashcards are blocked by a content hash. Cost is controlled by the input budgets, by
lower temperatures for structured tasks (0.1--0.3), by the per-plan credit reservation
in the v1 services (the free plan allows 50 credits per month; BYOK requests reserve
none), and by showing token usage per feature on the LLM Inspection page. Latency is
dominated by the provider; long tasks therefore stream progress or text to the
interface. With Gemini, the measured medians were 18.4~s to process a topic, 27.9~s for a Cornell
note, 52.9~s to generate a drill and 15.9~s for a Knowledge Hub answer (\Cref{tab:perf}); much of
this is waiting at the free-tier limit of 5 requests per minute.

\paragraph{Error handling.} Provider errors are mapped to typed errors: in the study
stack \texttt{classify\_error()} returns a \texttt{RecallError} with a code such as
\texttt{API\_KEY\_MISSING} or \texttt{LLM\_RATE\_LIMITED}, and in the v1 stack the
adapters raise \texttt{ProviderAuthError}, \texttt{ProviderRateLimitError},
\texttt{ProviderTimeoutError} or \texttt{ProviderUnavailableError}, which the routes turn
into HTTP errors with a clear message. Reserved credits are released when generation
fails.

\section{Code Structure and Standards Followed}\label{sec:code_structure}

\Cref{lst:tree} shows the folder structure of the repository (tracked files only, main
entries).

\begin{lstlisting}[caption={Repository folder structure (commit \texttt{60d40f9})},label={lst:tree}]
RecallAI/
|-- .github/workflows/     ci.yml, release.yml
|-- app/                   FastAPI backend (Python)
|   |-- main.py            app factory, middleware, legacy study routes
|   |-- api/               deps.py, middleware.py, v1/ (14 routers + router.py)
|   |-- core/              config, database (WAL), security, errors, migrations
|   |-- models/            repositories.py (29 repositories), schema_init.py
|   |-- schemas/           Pydantic request/response models
|   |-- services/          pipeline, chunker, embedding, retrieval, rag, quizzes,
|   |   |                  flashcards, reviews, scheduler (FSRS), analytics, ...
|   |   `-- ai/            AIService, adapters, base (v1 AI stack)
|   |-- llm_providers/     base, factory, gemini, openai, groq, ollama,
|   |                      rate_limiter, retry (study AI stack)
|   |-- extractors/        registry, pymupdf4llm, marker, marker_api plug-ins
|   |-- prompt_manager.py  PROMPT_REGISTRY (11 templates)
|   |-- socratic_drill.py  drill questions and grading
|   |-- llm_segment.py     segmentation, flashcards, chat
|   |-- database.py        per-user study database (14 tables)
|   |-- fsrs.py, toc_parser.py, pdf_extract.py, sync_service.py, ...
|   `-- tests/             22 pytest files (run in CI)
|-- desktop/               Tauri + React client
|   |-- src/
|   |   |-- api/           typed API client, auth, key storage
|   |   |-- components/    chat, dashboard, documents, flashcards, quizzes,
|   |   |                  review, settings, layout, ui, ...
|   |   |-- views/         page-level views (Dashboard, Documents, ...)
|   |   |-- store/         Redux Toolkit slices
|   |   `-- hooks/, contexts/, utils/, lib/
|   |-- src-tauri/         Rust shell, tauri.conf.json, sidecar binaries
|   `-- tests/e2e/         critical_user_journeys.mjs (Playwright)
|-- tests/                26 further pytest files and fixtures (not in CI)
|-- migrations/            PostgreSQL schema and composite indexes
|-- docs/database/         tbls schema documentation
|-- scripts/               build_sidecar.py, docker-entrypoint.sh
|-- Dockerfile, docker-compose.yml, requirements.txt
`-- run.bat, stop.bat, build.bat, AGENTS.md, README.md
\end{lstlisting}

\paragraph{Coding standards.}
\begin{itemize}[leftmargin=*]
  \item \textbf{Python.} The code follows PEP~8 naming and layout by
    convention~\cite{pep8}: \texttt{snake\_case} functions and modules,
    \texttt{PascalCase} classes, \texttt{UPPER\_CASE} constants, type hints on public
    functions, and docstrings on services and repositories.
  \item \textbf{TypeScript.} \texttt{tsconfig.app.json} enables \texttt{strict},
    \texttt{noFallthroughCasesInSwitch} and \texttt{noUncheckedSideEffectImports}
    (\texttt{noUnusedLocals} is off). Components use \texttt{PascalCase} file names,
    hooks start with \texttt{use}, and the API types are declared next to the client.
    \texttt{oxlint} runs in CI~\cite{oxlint} and reports 0 errors.
  \item \textbf{API conventions.} Plural nouns in paths, standard methods, a
    \texttt{\{"data": \ldots\}} envelope and one error format with fixed codes; every
    repository method of the v1 layer takes the workspace ID.
  \item \textbf{Security conventions.} Secrets come from the environment or generated
    local files, never from the code; user-visible keys are masked; uploads are checked
    by size and file signature.
  \item \textbf{Documentation.} Module docstrings describe each service; \texttt{README.md},
    \texttt{DEPLOYMENT.md}, \texttt{DESIGN.md} and \texttt{AGENTS.md} describe set-up,
    deployment, the visual design and the rules for AI coding agents (for example,
    test only with documents the user has uploaded).
\end{itemize}

\paragraph{Code review.}
The intended practice is review through pull requests on GitHub, with CI as an
automatic check on every pull request. In practice only one pull request (\#1, a
Copilot fix to the release workflow) was merged through GitHub; the rest of the work was
pushed directly to \texttt{main} (\Cref{sec:git}). Review was therefore informal: the
members checked each other's work in person at the weekly meetings before the lab
session. The Antigravity agent was also used to review the application and find bugs,
and its findings were checked by the team.
